% =====================================================================
% Section 12, Phase 2 research questions Q13–Q25 (humanised, single
% bundled figure per question from notebook_sections/)
% =====================================================================
\makesectioncover{Phase 2 \textbullet\ The Thirteen Questions}

\section{Phase 2 \textbullet\ Thirteen Research Questions}
\label{sec:phase2-questions}

The Phase 2 analysis notebook walks through thirteen questions, numbered
Q13 to Q25 to continue the Phase 1 numbering. Every question reads from
the cleaned and integrated outputs of Notebook 1 and from the relevant
Phase 1 CSVs. For each question we say what we are asking, where the
data came from, what method we used, and what the resulting visualisation
shows. The figure for every question is the bundled set of charts that
the analysis notebook produced for that question.

% --------------------------------------------------------------------
\subsection{Q13 \textbullet\ Metro coverage and density}

Did each metro line open into the population geography that existed at
the time of opening? An older station should not be judged against
2018 population alone, so we matched every metro station to its nearest
district and used the closest available population year (1996, 2006,
2017 or 2023).

Data: \fn{metro\_stations.csv} for the 89 stations with line and
opening year, joined to \fn{districts\_wide.csv} for district
centroids and yearly population. Method: a Spearman rank correlation
per line between opening year and adjacent-district population, plus
a scatter per line.

\figitem{phase2/sections/q13_notebook_visuals.png}{1.0}{Q13 bundled
visuals from the analysis notebook.}{fig:q13b}

What we see: Lines 1 and 2 (built before 2005) sit inside the older
dense core. Line 3, extended through 2024, has its newer Phase 3B and
3C stations pushing into denser districts like Imbaba and Mohandessin.
That part of the planning improved over time; the next question tests
whether it actually reached the 115 ghost terminals from Phase 1.

% --------------------------------------------------------------------
\subsection{Q14 \textbullet\ Ghost terminals near the new metro}

If Phase 1's ghost terminals (G1) sit close to a new Phase 3 metro
station, the terminal probably became redundant because riders shifted
to rail. If they sit far away, the structural problem from Phase 1 is
still there and the new rail did not address it.

Data: \fn{merged\_G\_underused\_terminals.csv} (the 115 ghosts) and
the post-2014 cohort of \fn{metro\_stations.csv}. Method: for each
ghost, take the haversine distance to the nearest post-2014 metro
station, keep the minimum, and bucket into 0--1 km, 1--2, 2--5, 5--10,
and \stat{$>$10 km}.

\figitem{phase2/sections/q14_notebook_visuals.png}{1.0}{Q14 bundled
visuals: distance buckets and the spatial diagnostic.}{fig:q14b}

The result: only \stat{9 of the 115 ghosts (8\%)} sit within 1 km of a
post-2014 metro station. \stat{97 of 115 (84\%)} are beyond 2 km. The
buckets read \stat{3 / 6 / 9 / 27 / 70} from nearest to farthest. The
new metro did not reach the underused-terminal cohort that Phase 1
flagged.

% --------------------------------------------------------------------
\subsection{Q15 \textbullet\ Metro and bus or microbus terminal integration}

Even where the metro and the older terminals do coexist, the question
is whether they integrate. A metro station near a high-route terminal
gives riders a quick transfer into the informal network; a station far
from any terminal is isolated infrastructure.

Data: metro stations with opening year plus the Phase 1 terminal
points. Method: for every metro station we found the nearest terminal,
attached its route count, passenger flow and ghost flag, and looked at
the per-line distribution of transfer distance.

\figitem{phase2/sections/q15_notebook_visuals.png}{1.0}{Q15 bundled
visuals from the analysis notebook.}{fig:q15b}

What it shows: older lines sit close to existing mawaqif and integrate
naturally. Some Line 3 stations, especially toward the eastern
extension, sit further from the informal terminal layer and therefore
need an explicit routing layer to be useful as a transfer point.

% --------------------------------------------------------------------
\subsection{Q16 \textbullet\ Fastest-growing districts and coverage}

A planning network that follows demand should serve growth as well as
standing population. Q16 takes the 15 districts with the highest CAGR
(2006 to 2017) and asks whether they have any new-mode coverage at
all.

Data: \fn{districts\_wide.csv} for CAGR plus the integrated station
layer for nearby station counts inside a 3 km buffer.

\figitem{phase2/sections/q16_notebook_visuals.png}{1.0}{Q16 bundled
visuals: top-15 fastest-growing districts and their coverage.}{fig:q16b}

The top-growth list is dominated by satellite cities such as New Cairo
1, 6th October City 1 and 3, Ash-Shuruq, and Sheikh Zayed. Established
dense districts like El Marg and Al-Warraq grow more slowly. Most of
the new metro extensions went to the satellite cities, which means a
lot of the slower-growth dense districts received little new coverage.

% --------------------------------------------------------------------
\subsection{Q17 \textbullet\ Population density vs. underserved score}

We test whether underservedness (Phase 1 G4) is just a density artefact
or a real coverage mismatch. If it were simply density, the score
would correlate trivially with population. If it correlates strongly
with population the underservedness is structural; those are the
districts where the gap is largest.

Data: \fn{q17\_dense\_underserved\_hexes.csv}, recomputed from Phase 1's
hex layer plus S6 district population. Method: a Spearman correlation,
a hex-bin joint distribution with marginal histograms, and a quadrant
overlay highlighting the dense and underserved cells.

\figitem{phase2/sections/q17_notebook_visuals.png}{1.0}{Q17 bundled
visuals.}{fig:q17b}

The result is \stat{$\rho = 0.548$, $p < 10^{-100}$}. The densest hexes
are also the most underserved. The market problem is not just density,
it is density combined with weak coverage.

% --------------------------------------------------------------------
\subsection{Q18 \textbullet\ Informal-transport share by density}

If informal microbus dependence concentrated only in the dense core,
it would be a poverty-pocket problem. If it sits flat across density
tiers, the informal network is the substrate everywhere.

Data: Phase 1 boarding (informal vs. total) joined to S6 district
population. Method: Spearman correlation plus an OLS slope.

\figitem{phase2/sections/q18_notebook_visuals.png}{1.0}{Q18 bundled
visuals: informal share by density.}{fig:q18b-tier}

The slope is statistically zero (\stat{$\rho = 0.025$, $p = 0.40$}).
Informal share averages around \stat{47\%} across every density tier.
Microbuses are the everyday substrate, not a poverty pocket.

% --------------------------------------------------------------------
\subsection{Q18b \textbullet\ Gap-closure matrix \textbullet\ the verdict}

Q18b cross-tabulates the four Phase 1 gaps (G1 to G4) against the
three post-2014 modes (Metro L3, LRT, BRT). Every cell is the
percentage of that gap's sites that sit within 2 km of a station of
that mode.

Data: the four gap-site sets (115 ghosts, 19 empty-return, around
1{,}340 vehicle-mismatch routes, 79 underserved hexes) buffered
against each mode's stations. Method: per cell, count gap sites
inside a 2 km haversine buffer of any station of that mode, report
as a percentage.

\figitem{phase2/sections/q18_notebook_visuals.png}{1.0}{The gap-closure
matrix is the most important single chart in this report. The bundled
Q18 figure shows it together with the per-tier informal share.}{fig:q18b}

The matrix verdict: \stat{no cell above 25\%}, most below \stat{16\%},
and the single best cell is BRT $\times$ G3 at \stat{25.0\%}. The
post-2014 build closes a quarter or less of any single gap.

% --------------------------------------------------------------------
\subsection{Q19 \textbullet\ GTFS coverage vs. Phase 1 reality}

There is a data gap as well as an infrastructure gap. We compare the
formal published GTFS network against the Phase 1 measured reality and
ask which Phase 1 stops sit outside an 800 m GTFS catchment.

Data: \fn{gtfs\_stops.csv} (S1) vs. Phase 1 boarding. Method: nearest
GTFS stop per Phase 1 stop with a coverage flag.

\figitem{phase2/sections/q19_notebook_visuals.png}{1.0}{Q19 bundled
visuals: GTFS coverage and the agency split.}{fig:q19b}

A meaningful share of informal-heavy stops sits beyond the 800 m
catchment. The data asymmetry is real and quantifiable.

% --------------------------------------------------------------------
\subsection{Q20 \textbullet\ BRT corridor vs. informal demand}

The BRT is the contrast case for the H3 hypothesis. Q20 ranks every
BRT station by the informal demand inside its 500 m buffer.

Data: \fn{brt\_stations.csv} (S5) plus Phase 1 boarding. Method: per
station, sum daily informal boardings within 500 m, rank descending.

\figitem{phase2/sections/q20_notebook_visuals.png}{1.0}{Q20 bundled
visuals: BRT stations ranked by nearby informal demand.}{fig:q20b}

A clear top tier of BRT stations carries thousands of informal
boardings within 500 m. That top tier is what drives the H3 result.

% --------------------------------------------------------------------
\subsection{Q21 \textbullet\ Fare per km, formal vs. informal}

Accessibility is not only spatial, it is also economic. Q21 asks how
much riders pay per km on each mode.

Data: \fn{cleaned\_routes.csv} (Phase 1) plus GTFS fare attributes.
Method: per route compute LE/km, group by mode, plot as a box.

\figitem{phase2/sections/q21_notebook_visuals.png}{1.0}{Q21 bundled
visuals: fare per kilometre by mode.}{fig:q21b}

Real numbers: formal bus is around \stat{0.22 LE/km}, formal metro
\stat{0.25 LE/km}, microbus \stat{0.62 LE/km} (about 3$\times$ formal),
tomnaya \stat{1.30 LE/km} (around 6$\times$ formal). The implication
is that for many origin-destination pairs the metro is unreachable, so
riders default to the higher per-kilometre microbus fare.

% --------------------------------------------------------------------
\subsection{Q22 \textbullet\ Metro under-performance vs. population}

Move from raw counts to a coverage expectation. We fit metro stations
as a function of district population. The residual (actual minus
predicted) is the under-coverage metric.

Data: \fn{metro\_stations.csv} plus S6 districts. Method: OLS
regression of stations on population (3 km radius per district
centroid), then rank residuals.

\figitem{phase2/sections/q22_notebook_visuals.png}{1.0}{Q22 bundled
visuals: top-15 worst residuals and the per-governorate distribution.}{fig:q22b}

The four worst residuals are \strval{Ad-Duqqi} ($-9.28$),
\strval{Sheikh Zayed} ($-9.09$), \strval{6th October City 1 \& 3}
($-8.52$) and \strval{Al-Hawamidiyah} ($-8.14$). All are dense or
fast-growing edges in Giza or Qalyubia. This is the ranking that the
final coverage-need map is built from.

% --------------------------------------------------------------------
\subsection{Q23 \textbullet\ Adly Mansour as a convergence node}

Adly Mansour is a symbolic interchange where Metro Line 3, the LRT,
the BRT and several Phase 1 informal terminals come together inside
2.5 km. We measure how its station density compares against 150 random
2.5 km clusters elsewhere in Cairo.

Data: integrated station and terminal layers plus Phase 1 boarding.
Method: count modes, count stations, then compute the percentile rank
of Adly Mansour against the random sample.

\figitem{phase2/sections/q23_notebook_visuals.png}{1.0}{Q23 bundled
visuals.}{fig:q23b}

\stat{Four modes within 2.5 km}, \stat{eight stations or terminals
total}, but only the \stat{21st percentile} of random Cairo clusters
by station density. Symbolically multimodal, statistically average.

% --------------------------------------------------------------------
\subsection{Q24 \textbullet\ K-Means coverage-need synthesis (AI Technique 2)}

Single-variable rankings cannot describe a market. The project needs
a segmentation, which is why Q24 runs K-Means on five engineered
features per district.

Data: S6 districts plus the integrated station layer. Method: K-Means
with \stat{$k=4$}, k-means++ initialisation, $n_{\text{init}}=50$,
features = \fn{log\_pop\_2023}, \fn{CAGR\,\%}, stations within 3 km,
informal stops within 3 km, and a Cairo-governorate indicator. We
re-ran across five seeds and got an Adjusted Rand Index of
\stat{1.00}, so the partition is fully stable. Silhouette at $k=4$ is
\stat{0.412}.

\figitem{phase2/sections/q24_notebook_visuals.png}{1.0}{Q24 bundled
visuals: elbow and silhouette, cluster heatmap, opportunity score,
parallel coordinates and the cluster radar.}{fig:q24b}

Four clusters come out:

\begin{itemize}
 \item Formal-Served Core: 31 districts, CAGR 0.19\%, about
 \stat{5.6M} residents.
 \item Peripheral Growth: 19 districts, CAGR 2.83\%, about
 \stat{9.2M} residents.
 \item Mixed (cluster 2): 12 districts, CAGR 2.17\%, about
 \stat{4.2M} residents.
 \item Low-Activity Outskirts: 6 districts, CAGR 11.93\%.
\end{itemize}

The first three together are \stat{62 districts and 18.9 million
residents}. That is the addressable market for any product that
bridges the two systems.

% --------------------------------------------------------------------
\subsection{Q25 \textbullet\ Masari bridge \textbullet\ informal to formal}

Q25 asks how connectable the informal layer is to the formal layer.
For every informal-heavy Phase 1 stop we find the nearest formal node
(GTFS, metro, LRT, BRT) and treat the connector as the missing
product layer.

\figitem{phase2/sections/q25_notebook_visuals.png}{1.0}{Q25 bundled
visuals: the Masari bridge schematic.}{fig:q25b}

Most informal-heavy stops sit within walking distance of some formal
node, so the bridges already exist physically. What is missing is the
software that surfaces them in one screen.
